Gli agenti basati su LLM possono affrontare compiti sempre più complessi, ma per collaborare tra loro, soprattutto quando appartengono a organizzazioni diverse, hanno bisogno di un protocollo comune. Il protocollo A2A risponde a questa esigenza, e il ciclo di vita dei Task ne è la parte centrale: è attraverso gli stati di un Task che un agente sa se il compito affidato a un altro è in corso, in attesa di informazioni, concluso o cancellato. Questa tesi si è occupata di come un agente produce correttamente questi stati, con particolare attenzione alla cancellazione.

Il lavoro è nato dall'esigenza, emersa durante il tirocinio presso HCL Software, di una gestione della cancellazione valida per qualsiasi agente creato dagli utenti, all'interno dell'SDK Python ufficiale e della versione 1.0 del protocollo. Il primo passo è stato l'aggiornamento del codice aziendale dalla versione 0.3 alla 1.0, mantenendo la compatibilità con i client esistenti. Il lavoro si è poi concentrato sulla classe \texttt{AgentExecutor}, e in particolare sui suoi due metodi, \texttt{execute} e \texttt{cancel}.

Per \texttt{execute} si è scelto di adottare un'unica strategia di esecuzione, che segue il grafo dell'agente sempre in streaming e pubblica un evento a ogni passo, lasciando al gestore delle richieste dell'SDK il compito di consegnare gli eventi al client nella modalità richiesta. In questo modo si evitano i due percorsi di codice paralleli della versione precedente, e la logica dell'agente non dipende più dal tipo di richiesta ricevuta. Per \texttt{cancel} è stato introdotto un registro dei task in esecuzione, che permette di distinguere un Task in corso da un Task in attesa di input. Le due situazioni vengono gestite in modo diverso, secondo il principio per cui lo stato di un Task è gestito da chi ne sta eseguendo l'elaborazione: se il Task è in esecuzione, \texttt{cancel} si limita a richiederne l'interruzione ed è \texttt{execute} a pubblicare lo stato \texttt{CANCELED}; se invece non c'è nessuna esecuzione attiva, è \texttt{cancel} a pubblicarlo. Entrambi i metodi sono raccolti in una classe comune, \texttt{BaseExecutor}, che può essere usata da agenti diversi senza modifiche.

I test svolti sui due agenti di esempio hanno confermato il comportamento atteso in tutti gli scenari considerati: l'esecuzione completa in modalità sincrona e in streaming, la cancellazione di un Task in esecuzione e di uno in attesa di input, la ripresa di un Task dopo una richiesta di input, il rifiuto delle richieste su Task inesistenti e delle cancellazioni ripetute, e l'esecuzione concorrente di più Task. La stessa classe ha gestito entrambi gli agenti senza modifiche.

Il lavoro ha anche mostrato che la cancellazione è un problema più delicato di quanto sembri. In \texttt{asyncio} la cancellazione è cooperativa, e l'interruzione avviene solo nei punti in cui il codice si sospende; inoltre una parte delle operazioni è svolta dall'SDK, e il risultato finale dipende da come i compiti sono divisi tra il framework e l'agente. Anche verificarla non è semplice: per inviare una richiesta di cancellazione serve l'identificativo del Task, generato dal server, e con Task di breve durata non c'è il tempo di leggerlo e usarlo. Per questo i test hanno richiesto agenti costruiti apposta, con una durata controllabile.

\section*{Sviluppi futuri}

Lo studio dell'SDK ha messo in luce alcuni casi che la soluzione attuale non copre, descritti nella Sezione~\ref{sec:limiti}, e che indicano i principali sviluppi futuri del lavoro.

Il primo riguarda le operazioni sui Task già terminati. Nella versione adottata dell'SDK il comportamento in questi casi non coincide sempre con quello previsto dalla specifica, e può dipendere dal momento in cui arriva la richiesta. Una soluzione è estendere il gestore delle richieste dell'SDK con una sottoclasse che controlli lo stato del Task prima di elaborare la richiesta; in alternativa, si può adottare una versione successiva dell'SDK che includa le correzioni già proposte nel suo repository.

Il secondo riguarda la cancellazione che arriva mentre \texttt{execute} sta pubblicando l'esito dell'elaborazione. In questa breve finestra l'esito già prodotto potrebbe andare perso ed essere sostituito dallo stato \texttt{CANCELED}. Il problema può essere affrontato con una variabile che impedisca di sovrascrivere un esito già prodotto, insieme a \texttt{asyncio.shield}, che protegge la pubblicazione dall'interruzione.

Il terzo riguarda la finestra che separa la fine di \texttt{execute} dall'aggiornamento dello stato memorizzato da parte dell'SDK. Questa finestra non si può chiudere dall'\texttt{AgentExecutor}, e richiederebbe un aggiornamento condizionato dello stato nello strato di persistenza del framework, che rifiuti qualsiasi transizione a partire da uno stato terminale.

Infine, la validazione potrebbe essere estesa con test automatici, che invochino direttamente il gestore delle richieste dell'SDK e controllino con precisione l'ordine degli eventi. Test di questo tipo renderebbero le verifiche ripetibili dopo ogni modifica del codice e permetterebbero di riprodurre anche i casi limite descritti, che con richieste inviate a mano non si possono ottenere. Sarebbero anche il presupposto per adottare la soluzione negli agenti reali dell'azienda, dove la gestione del ciclo di vita dei Task deve funzionare per qualsiasi agente creato dagli utenti.
